% Quarantine-only second-revision candidate proof.
% Date: 29 August 2026.
% This file is not a live-book input and has not been Rethlas-verified.
% It makes no novelty or priority claim.  The rejected first proof is
% and the rejected first revision are preserved under their original
% filenames.

\paragraph{Slope and orientation conventions.}
We use the type decomposition, branch space, residual curve, and canonical
graph metric of \cite[\S\S3.3--4.5]{Ducros24CurveStructure}.  The point
that was mishandled in the rejected first proof is the distinction between
two slopes.  If an analytic unit $f$ maps a sufficiently short section of
a branch $b$ by a dilation of canonical graph metrics with signed ratio
$\varepsilon r$, then
\begin{equation}
             \mathrm d\log_b|f|
                  =\varepsilon r [s(b):K]\in\mathbf Z
\label{eq:p213-normalized-versus-metric-slope}
\end{equation}
by \cite[\S4.5.21.4]{Ducros24CurveStructure}.  Thus the metric derivative
is $\varepsilon r$, not the integer in
\eqref{eq:p213-normalized-versus-metric-slope}.  All lattice and balancing
arguments below use the integer $\mathrm d\log_b|f|$; the canonical metric
is used only to identify the orientation of the image ray.

At a type--$2$ point $y$, the residual curve $C_y$ is projective, normal,
and integral over the actual degree-one residue field of $K$; it need not
be geometrically integral.  Since $y$ is interior, its branches are the
closed points of $C_y$ by
\cite[Theorem~4.2.10 and \S4.2.11.1]{Ducros24CurveStructure}.  A type--$3$
interior point has exactly two branches by
\cite[\S4.2.11.2]{Ducros24CurveStructure}.  For every analytic unit, the
normalized node relation
\begin{equation}
                    \sum_{b}\mathrm d\log_b|f|=0
\label{eq:p213-normalized-node-relation}
\end{equation}
holds on a boundaryless curve by
\cite[Theorem~4.5.22]{Ducros24CurveStructure}.

\begin{proof}[Proof of
\cref{lem:p213-variation-graph-source-cutoff}]
We first separate the active branches by normalized integer slope vectors.
Suppose that $y$ is of type~$2$, write $C=C_y$, and let
\[
                     b_1,\ldots,b_s
\]
be the branches on which $F_0$ is not locally constant.  There are only
finitely many of them: the variation graph of each coordinate of $F_0$ is
locally finite by \cite[\S\S4.5.21.1--4.5.21.2]
{Ducros24CurveStructure}.  Denote by $P_i\in C$ the closed point
corresponding to $b_i$.  If $s=0$, take $S$ to be empty and pass directly
to the source-star and cutoff construction below.  We henceforth assume
$s>0$ in the separation construction.

Weak approximation on $C$ gives rational functions
$r_1,\ldots,r_s\in\operatorname{Rat}(C)^\times$ such that
\begin{equation}
       \operatorname{ord}_{P_j}(r_i)=\delta_{ij}
       \qquad(1\leq i,j\leq s).
\label{eq:p213-variation-first-order-block-revised}
\end{equation}
For completeness, choose an effective divisor $D$ disjoint from all
$P_j$ and so positive that
\[
 H^1\!\left(C,
   \mathcal O_C\!\left(D-2P_i-\sum_{j\ne i}P_j\right)\right)=0.
\]
The resulting jet-evaluation map is surjective.  Prescribing a
uniformizer modulo $\mathfrak m_{P_i}^2$ at $P_i$ and the value $1$
modulo $\mathfrak m_{P_j}$ at every $P_j\ne P_i$ gives $r_i$.
This argument is over the actual constant field of $C$ and requires
neither geometric integrality nor a rational point.

Let $B$ be the finite set of closed points outside
$\{P_1,\ldots,P_s\}$ at which at least one $r_i$ has a zero or a pole.
A second jet-evaluation argument, with a fresh sufficiently positive pole
divisor disjoint from $B\cup\{P_1,\ldots,P_s\}$, gives
$h\in\operatorname{Rat}(C)^\times$ such that
\begin{equation}
 \operatorname{ord}_{P_i}(h)=0,
 \qquad
 \operatorname{ord}_{Q}(h)=1\quad(Q\in B).
\label{eq:p213-variation-collision-killer-revised}
\end{equation}
At any further zero or pole $R$ of $h$, all the $r_i$ have order zero by
the definition of $B$.

We next lift these rational functions and compute their correct slopes.
The field $\operatorname{Rat}(C)$ is the degree-one residue field of
$\mathcal H(y)$.  A norm-one representative in $\mathcal H(y)$ can be
approximated to error strictly less than one by an element of the dense
subfield $\kappa(y)$, without changing its residue.  Since $Y$ is normal
and $y$ is nonrigid,
\begin{equation}
                         \mathcal O_{Y,y}=\kappa(y)
\label{eq:p213-variation-normal-local-field-revised}
\end{equation}
by \cite[\S4.2.18]{Ducros24CurveStructure}.  Hence the approximants are
analytic unit germs.  Shrinking once, all of them are analytic units on a
common neighborhood of $y$.

For a branch $b_P$ corresponding to a closed point $P\in C$, choose a
norm-one lift $s_P$ of a residual uniformizer and orient $b_P$ away from
$y$.  Define
\begin{equation}
             d_P:=-\mathrm d\log_{b_P}|s_P|\in\mathbf Z_{>0}.
\label{eq:p213-variation-integral-branch-weight}
\end{equation}
The positivity follows from the fact that the residue of $s_P$ vanishes
at $P$, so $|s_P|$ decreases on a sufficiently short section of $b_P$.
The number $d_P$ is independent of the chosen norm-one uniformizer.  In
fact, \cite[\S4.2.11.1]{Ducros24CurveStructure} says that the refined
branch valuation is the composite of the original absolute value and the
discrete valuation at $P$.  If $q$ has secondary valuation zero, applying
the inequalities of \cite[\S4.2.6]{Ducros24CurveStructure} to $q$ and
$q^{-1}$ shows that $|q|$ is constant on a short branch section.  Therefore,
for every $g\in\operatorname{Rat}(C)^\times$ and every norm-one lift
$\widehat g$,
\begin{equation}
       -\mathrm d\log_{b_P}|\widehat g|
                 =d_P\operatorname{ord}_P(g).
\label{eq:p213-variation-residual-normalized-slope}
\end{equation}
The same argument for an arbitrary analytic unit germ $f$ gives
\[
                  -\mathrm d\log_{b_P}|f|\in d_P\mathbf Z.
\]

For a finite moment $G=(g_1,\ldots,g_m)$ that is nonconstant on $b$, set
\begin{equation}
 n_b(G):=\bigl(-\mathrm d\log_b|g_1|,\ldots,
                    -\mathrm d\log_b|g_m|\bigr)\in\mathbf Z^m,
 \qquad q_b(G):=\gcd_j|n_{b,j}(G)|.
\label{eq:p213-integral-slope-vector}
\end{equation}
By \eqref{eq:p213-normalized-versus-metric-slope}, the derivative of
$G_{\mathrm{trop}}$ with respect to canonical source length is
$n_b(G)/[s(b):K]$.  Consequently the primitive integral outgoing tangent
of the target ray is
\begin{equation}
                              u_b(G)=n_b(G)/q_b(G).
\label{eq:p213-primitive-integral-tangent}
\end{equation}
This is precisely where the factor $[s(b):K]$ cancels from the direction
but not from the metric dilation.

Let $S$ be the tuple of the lifts of $r_1,\ldots,r_s,h$.  On $b_i$, every
coordinate of $(F_0,S)$ has normalized slope divisible by $d_{P_i}$ and
the $r_i$-coordinate has slope exactly $d_{P_i}$.  Hence
$q_{b_i}(F_0,S)=d_{P_i}$ and the $S$-block of the primitive tangent is
$(e_i,0)$.  At a point $Q\in B$, the last coordinate in
\eqref{eq:p213-variation-collision-killer-revised} is one; at a further
zero or pole of $h$, the first $s$ coordinates vanish and the last does
not.  A branch activated only by $F_0$ has zero $S$-block.  It follows
that each selected ray is distinct from every other source ray.

If $y$ is of type~$3$, choose an element of $\mathcal H(y)^\times$ whose
modulus has non-torsion class modulo $|K^\times|$, approximate it by an
element of $\kappa(y)^\times$ without changing that modulus class, and use
\eqref{eq:p213-variation-normal-local-field-revised} to regard the result
as an analytic unit germ.  The description in
\cite[\S4.2.11.2]{Ducros24CurveStructure} shows that its modulus is
nonconstant on both branches, and
\eqref{eq:p213-normalized-node-relation} says that its two normalized
slopes are opposite.  Appending this one germ separates the two ray
germs.  We again call the appended block $S$.

Put $F=(F_0,S)$.  The union $\Gamma_F$ of the variation graphs of its
coordinates is closed and locally finite, every coordinate is locally
constant off $\Gamma_F$, and $F_{\mathrm{trop}}$ is piecewise a dilation
on $\Gamma_F$ by \cite[\S\S4.5.21.1--4.5.21.2]
{Ducros24CurveStructure}.  Shrink to an open tree $W$ about $y$ such that
$\Gamma_F\cap W$ is a finite star with its only vertex at $y$.  Thus every
active arm is mapped injectively and integral-affinely to a ray, whereas
$F_{\mathrm{trop}}$ is locally constant on every component of
$W\setminus\Gamma_F$.

The space $Y$ is locally holomorphically separated because it is good and
separated.  By \cite[Corollary~15.2.5]{CLD25}, there is a smooth function
$\psi$ with compact support in $W$ that is identically one near $y$.
Choose $\rho\in C^\infty(\mathbf R,[0,1])$ that is zero on a neighborhood
of $0$ and one on a neighborhood of $1$, and put $\chi=\rho\circ\psi$.
This is a smooth function: in every moment chart presenting $\psi$, it is
presented by the composite of the presenting smooth real function with
$\rho$.  Moreover,
\[
 0\leq\chi\leq1,\qquad
 \operatorname{supp}(\chi)\subseteq\operatorname{supp}(\psi)\Subset W,
 \qquad \chi=1\text{ near }y.
\]
Let
\[
        D:=\operatorname{supp}(\mathrm d'\chi)
             \cup\operatorname{supp}(\mathrm d''\chi).
\]
The compact set $D\cap\Gamma_F$ does not contain $y$.  Since every active
arm is mapped injectively away from $v$, one has
\[
                v\notin F_{\mathrm{trop}}(D\cap\Gamma_F).
\]
Choose an open neighborhood $O$ of $v$, small enough that it misses this
compact image and that
$F_{\mathrm{trop}}(\Gamma_F\cap\operatorname{supp}\chi)\cap O$ is the
finite separated fan germ described above.

Let now $T$ be as in the statement.  The support of $\chi T$ is compact,
so \cite[\S7.1.11]{CLD25} defines
$R_{T,F,\chi}=(F_{\mathrm{trop}})_*(\chi T)$ without any properness
assumption on $F_{\mathrm{trop}}$.  It is positive because $\chi\geq0$.
We prove directly that it is closed on $O$.  The Leibniz rule and the
closedness of $T$ give
\begin{equation}
 \mathrm d'R_{T,F,\chi}
     =(F_{\mathrm{trop}})_*(\mathrm d'\chi\wedge T),
 \qquad
 \mathrm d''R_{T,F,\chi}
     =(F_{\mathrm{trop}})_*(\mathrm d''\chi\wedge T).
\label{eq:p213-cutoff-differential}
\end{equation}
To pair either right-hand side with a target test form supported in $O$,
pull that form back by $F_{\mathrm{trop}}$.  On
$D\cap\Gamma_F$ its pullback has disjoint support by the choice of $O$.
On $D\setminus\Gamma_F$ the map $F_{\mathrm{trop}}$ is locally constant,
so the pullback of every positive-degree target form vanishes.  Both
right-hand sides of \eqref{eq:p213-cutoff-differential} therefore vanish
on $O$.

The same local-constancy observation shows that the support of the
pushforward in $O$ is contained in the displayed finite fan: a target
form supported away from the image of $\Gamma_F$ pulls back to zero.
Finally, the primitive $S$-blocks computed above show that a selected ray
germ has only its named source arm above it inside
$\Gamma_F\cap W$.  Its edge coefficient is therefore unmixed.  This
proves the lemma without a global target tube.
\end{proof}

\begin{proof}[Proof of
\cref{lem:p213-euclidean-one-dimensional-current-revised}]
The assertion is local on $\Omega$.  On a relatively compact coordinate
ball meeting only finitely many cells, positivity gives order zero by
\cite[Proposition~7.2.9]{CLD25}.  Thus $R$ is represented by a symmetric
positive-semidefinite matrix $M=(\mu_{ij})$ of Radon measures.  The
equation $\mathrm d'R=0$ says, column by column,
\begin{equation}
                         \sum_i\partial_i\mu_{ij}=0
                 \quad\text{in }\mathcal D'(\Omega).
\label{eq:p213-euclidean-column-divergence-revised}
\end{equation}

Consider a tube meeting the relative interior of one edge.  Choose linear
coordinates $(t,z_2,\ldots,z_N)$ in which the edge is $\{z=0\}$.  In
\eqref{eq:p213-euclidean-column-divergence-revised}, use a scalar test
function $z_a\psi(t,z)$, with arbitrary compactly supported $\psi$, for
each $a\geq2$.  Its tangential derivative vanishes on the edge and its
$a$-th normal derivative there is arbitrary.  Hence
$\mu_{aj}=0$ for $a\geq2$ and all $j$; symmetry gives
$\mu_{ia}=0$ as well.  The remaining equation is
$\partial_t\mu_{11}=0$.  A Radon measure on an interval with zero
distributional derivative is a constant multiple of Lebesgue measure,
and positivity makes the constant nonnegative.  In the integral affine
parametrization by the primitive tangent this is $b_e[e]$.

At a vertex $v$, subtract the edge currents on a finite star.  The
remainder is $A_v\delta_v$ for a positive-semidefinite matrix $A_v$.
The distributional column divergence of the full current at $v$ is, up to
one common sign,
\begin{equation}
 \left(\sum_{e\ni v}b_eu_{e,v}\right)\delta_v
       +\left(\sum_i(A_v)_{ij}\partial_i\delta_v\right)_{j=1}^N.
\label{eq:p213-euclidean-vertex-distribution-revised}
\end{equation}
Smooth test functions have independently prescribed value and gradient at
$v$.  First taking value zero and arbitrary gradient gives $A_v=0$;
then taking arbitrary value gives
\eqref{eq:p213-euclidean-vertex-balance-revised}.  Local finiteness gives
\eqref{eq:p213-euclidean-edge-decomposition-revised}.
\end{proof}

\begin{proof}[Proof of
\cref{prop:p213-normal-curve-arbitrary-valuation-current-revised}]
We first attach an intrinsic ratio to a source branch.  Let $y$ be of
type~$2$ or $3$, let $b$ be a branch at $y$, and choose a finite moment
$F$ supplied by
\cref{lem:p213-variation-graph-source-cutoff} that is nonconstant on $b$
and separates its ray germ.  Write
$n_b(F)$, $q_b(F)$, and $u_b(F)$ as in
\eqref{eq:p213-integral-slope-vector}--
\eqref{eq:p213-primitive-integral-tangent}.  By
\cref{lem:p213-euclidean-one-dimensional-current-revised}, the coefficient
of $(F_{\mathrm{trop}})_*(\chi T)$ on that ray is a number
$t_b(F)\geq0$.

Apply the same construction to the integration current and denote its
coefficient on the isolated primitive ray by
\begin{equation}
                              m_b(F)>0.
\label{eq:p213-actual-calibration-coefficient-revised}
\end{equation}
Strict positivity follows directly from the positive local degree in the
canonical calibration of
\cite[Proposition~14.3.1, \S14.3.2, and Definition~14.5.1]{CLD25},
because $F$ is nonconstant on $b$.  We deliberately do not identify
$m_b(F)$ with $q_b(F)$.  Over an inseparable constant-field extension,
the analytic calibration can retain a local-degree factor that disappears
from the normalized slope after completed-perfect-closure reduction.
For example, take $K=\mathbf F_p((t))$, $L=K(t^{1/p})$, and view
$\mathbf A^{1,\mathrm{an}}_L$ over $K$.  On a type--$3$ segment in
$\mathbf G_{m,L}^{\mathrm{an}}$, the coordinate moment has primitive
normalized slope $q_b(F)=1$, whereas the finite map to
$\mathbf A^{1,\mathrm{an}}_K$ has degree $p$ and hence
$m_b(F)=p$ by \cite[Remark~14.4.3]{CLD25}.  The quotient below retains
this factor.

The ratio
\begin{equation}
                              a_b:=t_b(F)/m_b(F)
\label{eq:p213-intrinsic-branch-ratio-revised}
\end{equation}
is independent of $F$ and of the cutoff.  Indeed, concatenate two moments
and, if necessary, append one separator block; use a common source cutoff
equal to one near $y$.  Write $H$ for this common refinement and let $A$
be the integral projection from its target to the target of $F$.  Since
$A n_b(H)=n_b(F)$, the primitive tangents satisfy
\begin{equation}
       A u_b(H)=r u_b(F),\qquad
       r=q_b(F)/q_b(H)\in\mathbf Z_{>0},\qquad
       A_*[u_b(H)]=r[u_b(F)].
\label{eq:p213-refinement-expansion-revised}
\end{equation}
The second identity follows by pulling back a $(1,1)$-form and changing
the integral edge parameter.  Functoriality of pushforward
\cite[\S7.1.11]{CLD25}, together with separation of the source ray, gives
\begin{equation}
                 t_b(F)=r\,t_b(H),\qquad
                 m_b(F)=r\,m_b(H).
\label{eq:p213-refinement-both-currents-revised}
\end{equation}
Thus the quotient in
\eqref{eq:p213-intrinsic-branch-ratio-revised} is unchanged.  Notice that
every inseparable or nilpotent local-degree factor is contained in both
occurrences of $m_b$ and is never discarded.
Replacing a cutoff by a smaller common cutoff does not change the ray-germ
coefficient.  Indeed, the difference of two such cutoffs vanishes near
$y$.  On the finitely many active arms of the common source star, the
image of its compact support is therefore bounded away from the target
vertex; off the variation graph every positive-degree target form pulls
back to zero.  The pushforward of the difference consequently vanishes on
a sufficiently small target germ.  This argument applies to both $T$ and
$\delta_Y$ and completes the proof of intrinsicity.

At a type--$3$ point, the two integral slope vectors are opposite by
\eqref{eq:p213-normalized-node-relation}.  Their greatest common divisors
are equal and their primitive tangents are opposite.  Applying
\eqref{eq:p213-euclidean-vertex-balance-revised} to the cutoff pushforward
first of $T$ and then of $\delta_Y$ gives equality, respectively, of the
two coefficients $t_b$ and of the two coefficients $m_b$.  Hence the two
branch ratios agree.

Let now $y$ be of type~$2$, put $C=C_y$, and retain the positive integers
$d_P$ from \eqref{eq:p213-variation-integral-branch-weight}.  For each
closed point $P\in C$, let $a_P$ be the ratio of its branch.  We also
attach an intrinsic positive calibration coefficient $m_P$ to that
branch.  Choose any separated moment that contains a norm-one lift of a
residual uniformizer at $P$.  Every coordinate has normalized slope
divisible by $d_P$, while the uniformizer coordinate has slope $d_P$.
Consequently the gcd of the moment's slope vector is $d_P$.  If two such
moments are replaced by their common refinement, its gcd is still $d_P$;
hence the expansion integer $r$ in
\eqref{eq:p213-refinement-expansion-revised} is one for both projections.
Equation \eqref{eq:p213-refinement-both-currents-revised} and cutoff
independence show that the integration-current coefficient is the same for
all these moments.  Denote this coefficient by $m_P>0$.  The corresponding
coefficient for $T$ is $a_Pm_P$.

Define two additive functionals on divisors by
\begin{equation}
       L_T([P])=a_Pm_P,
       \qquad
       L_Y([P])=m_P.
\label{eq:p213-divisor-functionals-revised}
\end{equation}
They are nonnegative on effective divisors.  We claim that both vanish on
principal divisors.

Fix $g\in\operatorname{Rat}(C)^\times$.  In one common moment append a
norm-one lift of $g$, a lift of a residual uniformizer at every point in
$\operatorname{Supp}(\operatorname{div}(g))$, and the finite separator
block constructed in the proof of
\cref{lem:p213-variation-graph-source-cutoff}.  On the branch at $P$, all
normalized slopes are divisible by $d_P$ and the uniformizer coordinate
has slope $d_P$.  Thus $q_{b_P}=d_P$, while the $g$-coordinate of the
primitive tangent is $\operatorname{ord}_P(g)$ by
\eqref{eq:p213-variation-residual-normalized-slope}.  Taking the
$g$-coordinate of
\eqref{eq:p213-euclidean-vertex-balance-revised} first for the cutoff
pushforward of $T$ and then for the cutoff pushforward of $\delta_Y$ gives
\begin{equation}
                         L_T(\operatorname{div}(g))=0.
\label{eq:p213-T-principal-zero-revised}
\end{equation}
\begin{equation}
                         L_Y(\operatorname{div}(g))=0.
\label{eq:p213-Y-principal-zero-revised}
\end{equation}

We use the following elementary Riemann--Roch consequence over the actual
constant field of $C$.  If
$L:\operatorname{Div}(C)\to\mathbf R$ is additive, nonnegative on
effective divisors, and zero on principal divisors, then $L=c_L\deg$.
Indeed, if $D$ has degree zero, choose an effective divisor $A$ with
\begin{equation}
                  \deg A+\chi(C,\mathcal O_C)>0.
\label{eq:p213-rr-positive-constant-revised}
\end{equation}
For every $n\in\mathbf Z$, Riemann--Roch gives
\[
 h^0\bigl(C,\mathcal O_C(A+nD)\bigr)
    \geq \deg A+\chi(C,\mathcal O_C)>0.
\]
Thus $A+nD$ is linearly equivalent to an effective divisor, and
$0\leq L(A)+nL(D)$ for every positive and negative integer $n$.
Consequently $L(D)=0$.  The functional factors through the subgroup
$\deg\operatorname{Div}(C)\subseteq\mathbf Z$, and positivity gives
$c_L\geq0$.

Apply this to \eqref{eq:p213-divisor-functionals-revised}, using
\eqref{eq:p213-T-principal-zero-revised} and
\eqref{eq:p213-Y-principal-zero-revised}.  There are constants
$c_T\geq0$ and $c_Y>0$ such that
\[
                  a_Pm_P=c_T\deg(P),
                  \qquad m_P=c_Y\deg(P).
\]
It follows that $a_P=c_T/c_Y$ for every branch at $y$.  This conclusion
uses the actual analytic calibration weights and therefore retains every
generic inseparable or nilpotent multiplicity.

We now identify the current locally.  Let $\alpha$ be an ordinary smooth
$(1,1)$-form with sufficiently small compact support near
$y\in Y^{[2,3]}$, and present it by a finite moment $F_0$ as in
\cite[Proposition~15.1.1(b)]{CLD25}.  Apply
\cref{lem:p213-variation-graph-source-cutoff}, taking the cutoff equal to
one near $\operatorname{supp}(\alpha)$ and shrinking the support of
$\alpha$, if necessary, into the inverse image of the target neighborhood
$O$.  In the local moment presentation, multiply the target representative
by a compactly supported target cutoff which is one on the compact image
of $\operatorname{supp}(\alpha)$.  This does not change its pullback.
On every branch on which $F_0$ varies, the preceding type--$2$ or
type--$3$ argument gives one common ratio $c_y$.  Edges activated only by
the separator block have zero $F_0$-tangent, so the target representative
vanishes on their primitive tangent.  The Euclidean lemma excludes an
additional vertex current.  The defining identity for compactly supported
pushforward therefore gives
\begin{equation}
                         \langle T,\alpha\rangle
                    =c_y\langle\delta_Y,\alpha\rangle.
\label{eq:p213-local-current-equality-revised}
\end{equation}

The locus $Y^{[2,3]}$ is nonempty and path-connected.  Indeed, a connected
normal analytic curve is an arcwise-connected graph by
\cite[Theorem~3.5.1 and \S1.3.5]{Ducros24CurveStructure}; it is not a
one-point graph because it is purely one-dimensional, and every interior
point of an open graph interval is of type $2$ or $3$.  Convexity of
$Y^{[2,3]}$ follows from \cite[\S4.5.11]{Ducros24CurveStructure}.
We spell out the propagation along an edge.  Fix a branch $b$ at a point
$y\in Y^{[2,3]}$.  The residual-uniformizer construction at type~$2$ and
the non-torsion-modulus construction at type~$3$ give an analytic unit
whose modulus varies on $b$.  After appending a separator moment and
shrinking, its variation graph on that arm has no vertex away from $y$,
and the tropical moment is injective and integral affine there.  Choose a
compact subsegment $J$ of the arm and a source cutoff that is one on a
neighborhood of $J$ and has compact support in a slightly longer segment
star.  The argument of
\cref{lem:p213-variation-graph-source-cutoff} shows that the pushforwards
of both $\chi T$ and $\chi\delta_Y$ are closed on a target neighborhood of
the image of $J$.  By
\cref{lem:p213-euclidean-one-dimensional-current-revised}, their
coefficients $t$ and $m$ are separately constant on this connected target
edge.  Their quotient $t/m$ is therefore constant along $J$.  At every
point of $J$ this quotient is the intrinsic ratio
\eqref{eq:p213-intrinsic-branch-ratio-revised}, because the same moment is
admissible there and that ratio is independent of the moment and cutoff.

Thus the ratio is locally constant along every open graph edge.  At a
type--$2$ or type--$3$ vertex the preceding star argument makes all
incident ratios equal.  Covering a path in the connected locus
$Y^{[2,3]}$ by finitely many such edge segments and vertex stars shows
that the constants $c_y$ in
\eqref{eq:p213-local-current-equality-revised} form one constant $c$ on
this locus.

Finally, every compactly supported ordinary $(1,1)$-form on the curve has
support in $Y^{[2,3]}$.  This follows from the local presentation
\cite[Proposition~15.1.1(b)]{CLD25} and the dimension argument in the
proof of \cite[Proposition~15.2.1(a)]{CLD25}.  A finite smooth partition
of unity on its compact support and
\eqref{eq:p213-local-current-equality-revised} give
$T=c\delta_Y$ globally.  By \cite[Example~8.6.4]{CLD25}, the integration
current is Bedford--Taylor, and the defining identity is stable under
nonnegative scalar multiplication.
\end{proof}

\begin{proof}[Proof of
\cref{prop:p213-curve-nonsmooth-locus-dichotomy-revised}]
The quasi-smooth locus is Zariski open by
\cite[\S3.1.11]{Ducros24CurveStructure}.  If it is nonempty, its
complement is a proper Zariski closed subset of the irreducible curve and
therefore has dimension zero.  On every strict affinoid domain this closed
analytic subspace is finite: its affinoid algebra is noetherian and
zero-dimensional, hence Artinian.  The strict affinoid Nullstellensatz
\cite[Theorem~4.5]{Tate71Rigid} makes all of its points rigid.  These local
finite descriptions show that the complement is locally finite in $Y$.
At every interior point, quasi-smoothness is equivalent to smoothness by
\cite[\S3.1.11.4]{Ducros24CurveStructure}; hence on a boundaryless curve
it is equivalent to rig-smoothness.  If the quasi-smooth locus is empty,
the second alternative holds.  The two alternatives are exclusive and
exhaustive.  This proves the assertion entirely on $Y$; no normalization
transfer is used.
\end{proof}

\begin{proof}[Verification of
\cref{ex:p213-nowhere-smooth-normal-curve-revised}]
Write $\alpha=t^{1/p}$, so
$L=K[\alpha]/(\alpha^p-t)$.  The algebraic $K$-curve
$\operatorname{Spec}L[T]$ is integral and regular, and its analytification
is good, separated, strict, and boundaryless.  It is normal and purely
one-dimensional.  In characteristic $p$, differentiating
$\alpha^p-t=0$ gives no relation on $\mathrm d\alpha$, and therefore
\begin{equation}
 \Omega^1_{L[T]/K}
      \simeq L[T],\mathrm d\alpha
           \oplus L[T],\mathrm dT.
\label{eq:p213-nowhere-smooth-differentials-revised}
\end{equation}
At every analytic point the fibre of the differential sheaf has dimension
two over the completed residue field, whereas the local analytic dimension
is one.  The criterion in \cite[\S3.1.11]{Ducros24CurveStructure} shows
that the curve is nowhere quasi-smooth over $K$.  It contains nonrigid
Gauss points, so its nonsmooth locus is not a locally finite rigid set.
\end{proof}
